\section{Identification by finite-matrix Dyson Brownian motion}
\label{sec:dbm}

Proposition~\ref{prop:gaussian-insertion} approximates every microscopic
subsequential law by the eigenvalues of a finite diagonal matrix plus GOE
noise. We now apply the fixed-energy theorem of Landon, Sosoe, and
Yau~\cite{LSY2019} to these finite matrices.
Two details require attention: their limiting density must be normalized
correctly, and the correlation-function conclusion must be converted into
the Laplace-functional statement of Theorem~\ref{thm:main}.

\subsection{The fixed-energy input}

For a real diagonal matrix \(V=\diag(v_1,\ldots,v_m)\), write
\[
 m_V(\zeta)=\frac1m\sum_{j=1}^m\frac1{v_j-\zeta}.
\]
For positive parameters \(g,G\), we call \(V\) \((g,G)\)-regular
if, for fixed positive constants
\(c,C,C_V\),
\begin{equation}
 \label{eq:dbm-regularity}
 c\le \Im m_V(e+i\eta)\le C
 \quad (|e|\le G,\ g\le\eta\le10),
 \qquad \|V\|\le m^{C_V}.
\end{equation}
The parameters may depend on \(m\), and satisfy, for a fixed
\(\delta>0\),
\begin{equation}
 \label{eq:dbm-parameter-domain}
 m^{\delta-1}\le g\le m^{-\delta},
 \qquad G\le m^{-\delta}.
\end{equation}
Let \(W_m\) be standard normalized GOE, with entry variances \(1/m\)
off the diagonal and \(2/m\) on the diagonal. The deterministic comparison
density for \(V+\sqrt t\,W_m\) is defined by the free-convolution transform
\begin{equation}
 \label{eq:dbm-free-convolution}
 m_{\mathrm{fc},t}(\zeta)
   =m_V\bigl(\zeta+t\,m_{\mathrm{fc},t}(\zeta)\bigr),
 \qquad
 \rho_{\mathrm{fc},t}(e)
   =\frac1\pi\lim_{\eta\downarrow0}
          \Im m_{\mathrm{fc},t}(e+i\eta).
\end{equation}
Here the solution with positive imaginary part is used in the upper
half-plane.

We state the specialization of the fixed-energy theorem that we need.
Its density convention is important. If \(p_H^{(m)}\) denotes the
symmetrized joint eigenvalue density of a random matrix \(H\), normalized
to have integral one, then
\[
 p_H^{(a)}(y_1,\ldots,y_a)
   =\int_{\RR^{m-a}}p_H^{(m)}(y_1,\ldots,y_m)
        \,d y_{a+1}\cdots\,d y_m
\]
also has integral one. It is not the factorial moment density.

\begin{theorem}[Fixed-energy DBM input]
\label{thm:dbm-input}
Suppose a sequence of deterministic diagonal matrices \(V\) of dimension
\(m\to\infty\) satisfies \eqref{eq:dbm-regularity} and
\eqref{eq:dbm-parameter-domain}, with fixed constants. Fix \(\sigma>0\),
and suppose
\begin{equation}
 \label{eq:dbm-time-window}
 g m^\sigma\le t\le m^{-\sigma}G^2.
\end{equation}
Assume also that \(\rho_{\mathrm{fc},t}(0)=1/\pi\).
For every fixed positive integer \(a\) and
\(O\in C_c^\infty(\RR^a)\), the difference
\begin{align}
 \label{eq:dbm-marginal-test}
 &\int_{\RR^a}O(\boldsymbol\alpha)
   p_{V+\sqrt t\,W_m}^{(a)}
      \left(\frac{\pi\alpha_1}{m},\ldots,
            \frac{\pi\alpha_a}{m}\right)
       \,d\boldsymbol\alpha \notag\\
 &\hspace{12mm}-
 \int_{\RR^a}O(\boldsymbol\alpha)
   p_{W_m}^{(a)}
      \left(\frac{\pi\alpha_1}{m},\ldots,
            \frac{\pi\alpha_a}{m}\right)
       \,d\boldsymbol\alpha
 \longrightarrow0.
\end{align}
Consequently, the expected sums of \(O\) over ordered tuples of
distinct points of the two eigenvalue processes, both rescaled by
\(m/\pi\), differ by \(o(1)\).
\end{theorem}

\begin{proof}
The marginal statement is \cite[Theorem~2.2]{LSY2019}, with
Definition~2.1 and Equations~(2.5)--(2.9) there, specialized to energy zero and equal densities
\(\rho_{\mathrm{fc},t}(0)=\rho_{\mathrm{sc}}(0)=1/\pi\).
The energy lies in the required interior window for any fixed interior
fraction. That theorem gives a power-decaying bound for
\eqref{eq:dbm-marginal-test}. To pass to the last assertion, multiply
each integral by
\[
 \frac{(m)_a}{(m/\pi)^a},
 \qquad
 (m)_a=m(m-1)\cdots(m-a+1).
\]
This is the Jacobian and ordered-tuple multiplicity for the unit-mass
marginal convention, and it tends to \(\pi^a\) for fixed \(a\).
\end{proof}

Theorem~\ref{thm:dbm-input} applies to arbitrary deterministic
real diagonal matrices satisfying its displayed hypotheses.
The same specialization is recorded in \cite[Theorem~6.1]{bu}.

\subsection{Regularity and the density adjustment}

Fix a microscopic subsequential law \(\nu\) supplied by
Proposition~\ref{prop:marked-limits}. We use the parameters and
conditional matrices of Proposition~\ref{prop:gaussian-insertion}:
\[
 K=r^{1/100},\qquad T=r^\tau,\qquad
 0<\tau<\frac{\chi}{100000},\qquad
 I=\{j:|x_j|\le K\},\qquad m_I=|I|,
\]
and
\begin{equation}
 \label{eq:dbm-raw-matrix}
 H_r=D_I+s_1\Id+\sqrt{s_2}\,\mathcal W_r,
 \qquad |s_1|\le C T,\qquad cT\le s_2\le CT.
\end{equation}
Conditional on \(\mathcal E\), \(\mathcal W_r\) is raw GOE of dimension
\(m_I\), with variances \(1\) and \(2\). All following almost-sure
assertions concern the conditioning input \(\mathcal E\) defined in
Section~\ref{sec:insertion}. In particular,
the discrepancy bound of Lemma~\ref{lem:discrepancy} gives
\begin{equation}
 \label{eq:dbm-dimension}
 m_I=2\rho K+O(K^{1-\chi}),
 \qquad
 T=m_I^{100\tau+o(1)}.
\end{equation}
Thus \(m_I>0\) eventually. We may decrease \(\chi\) so that
\(\chi<1/4\).

Put \(\alpha_*=\pi\rho\). Scaling \eqref{eq:dbm-raw-matrix} by
\(\alpha_*/m_I\) gives the normalized GOE-addition model
\begin{equation}
 \label{eq:dbm-normalized-matrix}
 \frac{\alpha_*}{m_I}H_r
   \ \stackrel{\mathrm{law}}{=}\
       V_*+\sqrt{t_*}\,W_{m_I},
 \qquad
 V_*=\frac{\alpha_*}{m_I}(D_I+s_1\Id),
 \qquad
 t_*=\frac{\alpha_*^2s_2}{m_I}.
\end{equation}

The following verification adapts the calculation in
\cite[Section~6.2]{bu}. We give it in full for the random coefficients
constructed here.

\begin{lemma}
\label{lem:density-adjustment}
For almost every realization of \(\mathcal E\), the matrices \(V_*\) are eventually regular
with
\[
 g=\frac{T^{1/2}}{m_I},\qquad G=m_I^{-1/8}.
\]
They satisfy the hypotheses
\eqref{eq:dbm-parameter-domain} and \eqref{eq:dbm-time-window} with
\(m=m_I\) and \(\delta=\sigma=\tau\). Moreover,
\[
 \rho_{\mathrm{fc},t_*}(0)\longrightarrow\frac1\pi.
\]
If \(c_r=\pi\rho_{\mathrm{fc},t_*}(0)\), then eventually \(c_r>0\),
and the data \(c_rV_*\) and time \(c_r^2t_*\) satisfy
Theorem~\ref{thm:dbm-input}, including its exact density condition.
\end{lemma}

\begin{proof}
First, \(\|V_*\|\) is bounded by
\eqref{eq:dbm-dimension}, since \(T=o(K)\).
We verify a slightly enlarged regularity domain,
\[
 |\Re\zeta|\le2G,\qquad g/2\le\Im\zeta\le20,
\]
so that subsequent dilation by \(c_r=1+o(1)\) will be harmless.
The normalized resolvent trace is
\[
 m_{V_*}(\zeta)=
   \alpha_*^{-1}\sum_{j\in I}
        \frac1{x_j-(m_I\zeta/\alpha_*-s_1)}.
\]
Writing \(e=m_I\Re\zeta/\alpha_*-s_1\) and
\(u=m_I\Im\zeta/\alpha_*\), the enlarged domain has
\begin{equation}
 \label{eq:dbm-original-domain}
 |e|=O(K^{7/8}+T)=o(K),\qquad
 cT^{1/2}\le u\le CK.
\end{equation}
The lower height exceeds a fixed positive power of \(K\).
Proposition~\ref{prop:transform-regularity}, with a radius comparable to \(K\)
and any fixed exponent smaller than \(50\tau\), therefore gives
\[
 \sum_j\frac{u}{(x_j-e)^2+u^2}
       =\pi\rho+o(1)
\]
uniformly in \eqref{eq:dbm-original-domain}.
For \(u/K\) sufficiently small, the omitted terms with \(|x_j|>K\)
sum to \(O(u/K)\), uniformly for \(|e|=o(K)\), by the linear count
bound from Lemma~\ref{lem:discrepancy}. Hence the retained sum
has a positive lower bound. For \(u/K\) bounded below, every retained
point has \(|x_j-e|\le2K\) eventually, so
\[
 \sum_{j\in I}\frac{u}{(x_j-e)^2+u^2}
   \ge \frac{m_Iu}{4K^2+u^2}\ge c>0.
\]
Its upper bound follows by positivity from the full sum.
This proves regularity on the enlarged domain.

With \(m=m_I\), the relevant powers are
\[
 g=m^{-1+50\tau+o(1)},\qquad
 t_*=m^{-1+100\tau+o(1)},\qquad G=m^{-1/8}.
\]
In particular,
\[
 \frac{t_*}{g m^\tau}=m^{49\tau+o(1)}\longrightarrow\infty,
 \qquad
 \frac{t_*}{m^{-\tau}G^2}
     =m^{-3/4+101\tau+o(1)}\longrightarrow0.
\]
The inequalities \(m^{\tau-1}\le g\le m^{-\tau}\) and
\(G\le m^{-\tau}\) follow as well. Our bound on \(\tau\) makes all
these exponent inequalities strict.

It remains to determine the free-convolution density. At \(\zeta=0\),
Equation~\eqref{eq:dbm-free-convolution} becomes
\begin{equation}
 \label{eq:dbm-boundary-equation}
 \mathfrak m=F_r(\mathfrak m),\qquad
 F_r(\mathfrak m)=\alpha_*^{-1}\sum_{j\in I}
     \frac1{x_j-(\alpha_*s_2\mathfrak m-s_1)}.
\end{equation}
For \(|\mathfrak m-i|\le1/4\), the argument
\(z=\alpha_*s_2\mathfrak m-s_1\) has \(\Im z\asymp T\) and \(|z|=O(T)\).
Propositions~\ref{prop:transform-regularity} and
\ref{prop:marked-resolvent} imply
\[
 \PV\sum_j\frac1{x_j-z}=i\pi\rho+o(1)
\]
uniformly in this disk, with the symmetric principal-value convention
of Section~\ref{sec:resolvents}. Furthermore,
\[
 \PV\sum_{|x_j|>K}\frac1{x_j-z}
   =\PV\sum_{|x_j|>K}\frac1{x_j}
      +\sum_{|x_j|>K}\frac{z}{x_j(x_j-z)}
   =O(K^{-\chi}+T/K).
\]
The first estimate follows by partial summation in
Lemma~\ref{lem:discrepancy}; the second uses its linear count bound and
\(|z|=o(K)\). Thus \(F_r(\mathfrak m)=i+o(1)\) uniformly on
\(|\mathfrak m-i|\le1/4\).

Rouch\'e's theorem applied to \(\mathfrak m-F_r(\mathfrak m)\) and \(\mathfrak m-i\) shows that
\eqref{eq:dbm-boundary-equation} has exactly one zero, counted with
multiplicity, in this disk for all large \(r\). The zero \(\mathfrak m_r\) is
therefore simple, and \(\mathfrak m_r\to i\).
The implicit function theorem continues it to a solution near
\(\zeta=0\) of the full fixed-point equation, still with positive
imaginary part. This solution is the physical branch. Indeed, every
upper-half-plane solution \(\mathfrak m\), at \(\Im\zeta>0\), satisfies
\begin{equation}
 \label{eq:dbm-uniqueness-identity}
 \frac1{m_I}\sum_{j\in I}
   |(V_*)_{jj}-\zeta-t_*\mathfrak m|^{-2}
       =\frac{\Im \mathfrak m}{\Im\zeta+t_*\Im \mathfrak m}<\frac1{t_*}.
\end{equation}
If two distinct such solutions existed, subtracting their equations
and dividing by their difference would give
\[
 1=\frac{t_*}{m_I}\sum_{j\in I}
  \frac1{((V_*)_{jj}-\zeta-t_*\mathfrak m_1)
          ((V_*)_{jj}-\zeta-t_*\mathfrak m_2)}.
\]
Cauchy--Schwarz and \eqref{eq:dbm-uniqueness-identity} make the
absolute value of the right-hand side strictly less than one.
This contradiction proves uniqueness and identifies the analytic
continuation. Its boundary value gives
\(\rho_{\mathrm{fc},t_*}(0)=\Im \mathfrak m_r/\pi\to1/\pi\).

Finally, simultaneous replacement of \(V_*,t_*\) by
\(c_rV_*,c_r^2t_*\) dilates the free-convolution law by \(c_r\).
For example its transform is
\(c_r^{-1}m_{\mathrm{fc},t_*}(\zeta/c_r)\), as follows directly
from \eqref{eq:dbm-free-convolution}. Its density at zero is exactly
\[
 c_r^{-1}\rho_{\mathrm{fc},t_*}(0)=\frac1\pi.
\]
Since \(c_r\to1\), the enlarged regularity domain proves
\eqref{eq:dbm-regularity} for \(c_rV_*\) with the original \(g,G\).
The strict power margins preserve \eqref{eq:dbm-time-window} for
\(c_r^2t_*\), completing the proof.
\end{proof}

The conclusion holds for almost every realization of \(\mathcal E\), which is
enough for applying Theorem~\ref{thm:dbm-input} to a deterministic
sequence. No assertion of switching invariance conditional on \(\nu\)
is used. The rescaled eigenvalues supplied by that theorem are, by
\eqref{eq:dbm-normalized-matrix}, precisely
\begin{equation}
 \label{eq:dbm-final-coordinates}
 \frac{m_I}{\pi}\frac{c_r\alpha_*}{m_I}
       \lambda_j(H_r)
       =c_r\rho\,\lambda_j(H_r).
\end{equation}
Define \(c_r=1\) at any exceptional initial values where the preceding
construction is not defined; these values have no effect on the limit.

\subsection{The GOE limit and uniform count moments}

Let
\[
 \mathcal G_m=\sum_{j=1}^m
       \delta_{(m/\pi)\lambda_j(W_m)}.
\]
These point processes converge vaguely in law, along all dimensions,
to a common locally finite point process, which we denote by
\(\mathcal S_{\mathrm{GOE}}\). To check normalization directly,
the eigenvalues of \(\sqrt m\,W_m\) have joint density proportional to
\[
 \exp\left(-\frac14\sum_j y_j^2\right)
       \prod_{i<j}|y_i-y_j|.
\]
This follows from orthogonal diagonalization of the matrix density
\(\exp(-\Tr H^2/4)\); the change in an off-diagonal coordinate under
an infinitesimal basis rotation is the corresponding eigenvalue
difference, yielding the Vandermonde Jacobian.
Theorem~1 of Valk\'o and Vir\'ag~\cite{ValkoVirag2009}, for \(\beta=1\) and center zero,
gives vague convergence after multiplying these eigenvalues by
\(2\sqrt m\). Its bulk condition
\(m^{1/6}(2\sqrt m-0)\to\infty\) holds. Dividing the resulting
coordinates by \(2\pi\) gives \(\mathcal G_m\), as claimed.

Weak convergence alone does not justify termwise expectations of
correlation expansions. The following bound makes the finite inclusion--exclusion brackets
converge uniformly for the reference ensemble.

\begin{lemma}
\label{lem:goe-counts}
For every compact interval \(J\subset\RR\) and every \(b>0\),
\[
 \sup_{m\ge1}\EE\exp\bigl(b\,\mathcal G_m(J)\bigr)<\infty.
\]
\end{lemma}

The proof is given in Appendix~\ref{app:goe-counts}.  It uses Gaussian
decimation to compare the GOE count with a determinantal count, followed
by a uniform bound on its projection kernel.


\subsection{From correlation functions to Laplace functionals}

We use the finite Bonferroni comparison in \cite[Section~6.4]{bu}
to pass from conditional DBM convergence to the unconditional
point-process law.
Let \(\nu^{(r)}=\sum_{j=1}^{m_I}\delta_{\lambda_j(H_r)}\) denote
the raw eigenvalue measure in \eqref{eq:dbm-raw-matrix}.
Proposition~\ref{prop:gaussian-insertion} asserts
\begin{equation}
 \label{eq:dbm-insertion-limit}
 \nu^{(r)}\ \Longrightarrow\ \nu
\end{equation}
in local vague topology, unconditionally. We will identify the law
of \(\rho\)-dilated \(\nu\) by first applying DBM conditional on \(\mathcal E\).

Fix a nonnegative \(h\in C_c^\infty(\RR)\), and put
\(f=1-e^{-h}\), so \(0\le f\le1\). For a finite point measure
\(\zeta=\sum_i\delta_{y_i}\), define
\[
 A_0(\zeta)=1,\qquad
 A_a(\zeta)=\frac1{a!}
    \sum_{i_1,\ldots,i_a\ {\rm distinct}}
        \prod_{j=1}^a f(y_{i_j})\quad(a\ge1),
 \qquad
 B_\ell(\zeta)=\sum_{a=0}^{\ell}(-1)^aA_a(\zeta).
\]
The finite inclusion-exclusion inequalities give
\begin{equation}
 \label{eq:dbm-bonferroni}
 B_{2\ell+1}(\zeta)
   \le \exp\left(-\int h\,d\zeta\right)
   \le B_{2\ell}(\zeta).
\end{equation}
One can see these inequalities by taking independent Bernoulli
variables of success probabilities \(f(y_i)\): the middle quantity
is their probability of no successes, and the displayed sums are
the Bonferroni bounds for that event.

Conditional on almost every realization of \(\mathcal E\), let \(\zeta_r\) have the point
coordinates \(c_r\rho\,\lambda_j(H_r)\).
Lemma~\ref{lem:density-adjustment} and Theorem~\ref{thm:dbm-input},
applied with the test \(\prod_{j=1}^a f(\alpha_j)\), imply for each
fixed \(\ell\) that
\[
 \EE[B_\ell(\zeta_r)\mid\mathcal E]
       -\EE B_\ell(\mathcal G_{m_I})\longrightarrow0.
\]
The expectation in the second term is over a reference GOE matrix
of the now deterministic conditional dimension \(m_I\).
To close the brackets uniformly in that dimension, choose a compact
interval \(J\) containing \(\supp h\). Then
\[
 A_a(\mathcal G_m)\le
     \binom{\mathcal G_m(J)}a
     \le 2^{-a}3^{\mathcal G_m(J)}.
\]
The last inequality follows from
\(\sum_{a\ge0}2^a\binom{N}{a}=3^N\).
Lemma~\ref{lem:goe-counts} consequently gives
\begin{equation}
 \label{eq:dbm-bracket-width}
 \sup_m\EE\bigl[B_{2\ell}(\mathcal G_m)
                 -B_{2\ell+1}(\mathcal G_m)\bigr]
       \le C_J2^{-(2\ell+1)}.
\end{equation}
Using \eqref{eq:dbm-bonferroni} for both models, first sending
\(r\to\infty\) at fixed \(\ell\), and then sending
\(\ell\to\infty\), proves
\begin{equation}
 \label{eq:dbm-conditional-laplace}
 \EE\left[\exp\left(-\int h\,d\zeta_r\right)\,\middle|\,\mathcal E\right]
   \longrightarrow
 \EE\exp\left(-\int h\,d\mathcal S_{\mathrm{GOE}}\right).
\end{equation}
Here we also used the vague GOE limit and \(m_I\to\infty\).
Only the reference count moments were needed in
\eqref{eq:dbm-bracket-width}; no uniform moment bound for the
deformed model has been assumed. Since both sides of the conditional
Laplace comparison lie in \([0,1]\), bounded convergence now removes
the conditioning in \eqref{eq:dbm-conditional-laplace}.

For \(c>0\), let \(\mathcal D_c\zeta\) denote the point measure
obtained by multiplying all positions of \(\zeta\) by \(c\).
The map \((c,\zeta)\mapsto\mathcal D_c\zeta\) is continuous when
\(c\) tends to a positive number and \(\zeta\) converges locally
vaguely. Indeed, the inverse images of a fixed compact support
remain in a common compact interval, and the corresponding test
functions converge uniformly there; vague convergence also bounds
the masses on a slightly larger compact interval.
By Lemma~\ref{lem:density-adjustment}, \(c_r\to1\) almost surely.
Combining this with \eqref{eq:dbm-insertion-limit} and the usual
joint convergence with a constant yields
\[
 \zeta_r=\mathcal D_{c_r\rho}\nu^{(r)}
       \ \Longrightarrow\ \mathcal D_\rho\nu.
\]
The Laplace functional of a compactly supported continuous
nonnegative test is bounded and continuous in local vague topology.
Thus \eqref{eq:dbm-conditional-laplace}, after averaging, identifies
\begin{equation}
 \label{eq:dbm-limit-identification}
 \EE\exp\left(-\int h\,d\mathcal D_\rho\nu\right)
   =\EE\exp\left(-\int h\,d\mathcal S_{\mathrm{GOE}}\right)
 \qquad (0\le h\in C_c^\infty(\RR)).
\end{equation}
For any \(0\le h\in C_c(\RR)\), convolution with a nonnegative
smooth approximate identity gives nonnegative smooth functions
converging uniformly to \(h\), all supported in one fixed compact
interval. Both limiting point measures have finitely many points
there almost surely. The corresponding Laplace variables converge
pointwise and are bounded by one, so bounded convergence extends
\eqref{eq:dbm-limit-identification} to every such \(h\).

\begin{proof}[Proof of Theorem~\ref{thm:main}]
Choose the fixed compact bulk interval \(J\) with
\(I_{d,\kappa}\subset\operatorname{int}J\).
The finitely many smallness choices in Proposition~\ref{prop:tree-stability}
and the subsequent arguments give one positive number \(w_0(d,\kappa)\).
Fix any \(0<w<w_0(d,\kappa)\) and \(E\in I_{d,\kappa}\).
Corollary~\ref{cor:tree-density} gives a continuous density on an open
neighborhood of \(I_{d,\kappa}\), strictly positive there; write
\(\rho=\rho_{d,w}(E)>0\).

Start with any sequence of admissible graph sizes tending to
infinity. Proposition~\ref{prop:marked-limits} supplies a further
subsequence on which the unscaled microscopic measures
\(\nu_n\) converge in law to a process \(\nu\) of the form used above.
The graph point measure in the theorem is
\(\mathcal D_\rho\nu_n\). Bounded continuity of its Laplace
functional and \eqref{eq:dbm-limit-identification} show that its
Laplace expectation, along this further subsequence, tends to the
Laplace expectation of \(\mathcal S_{\mathrm{GOE}}\).
The reference GOE expectations at dimension \(n\) tend to the same
value, along all sizes. Every original subsequence therefore has a
further subsequence along which their difference tends to zero,
which proves the required full-sequence limit. All spectral
estimates and negligible exclusions used to construct the
subsequential law were events in the original simple-graph model;
no conditioning on connectedness, bipartiteness, or any additional
graph property enters the conclusion. The constants in estimates that
use spectral averaging may depend on this fixed positive \(w\).
The argument therefore proves the assertion for every fixed
\(w\in(0,w_0(d,\kappa))\), with the limit in graph size taken afterward.
\end{proof}
